\chapter{Glossary}
\label{app:glossary}

\begin{description}[leftmargin=!,labelwidth=32mm,style=nextline,font=\normalfont\bfseries]

\item[Bond force] The downward load the tool holds against the pad while
  welding, produced by the force coil and set in grams.

\item[Configuration] A named set of the 31 bonding parameters --- a recipe.
  Configurations belong to one bonding mode; the machine stores up to 32.

\item[Constant force] The light holding force carried during descent and
  between process steps. Enough to keep the bond lever controlled, not enough
  to weld.

\item[Contact sensor] The electrical connection between tool and workpiece.
  The machine uses it to detect touchdown on the way down, and wire release on
  the way up.

\item[Energy] The total ultrasonic energy delivered into a weld, in
  millijoules --- the \emph{dose}. Together with power it determines how long
  the weld takes.

\item[Force coil] The voice-coil actuator that produces bond force. Its
  accuracy depends on the force calibration; see Chapter~\ref{ch:calibration}.

\item[Homing] The Y axis's search for its limit switch at power-up,
  establishing where zero is. The machine will not move until this succeeds.

\item[Kink] The sharp bend formed at the heel of the first bond as the tool
  rises, which starts the loop. Set by \menuitem{Kink Height}.

\item[Loop] The arc of wire between the two bonds. Its apex is set by
  \menuitem{Loop Height}.

\item[LVDT] The sensor that measures Z height. It reads absolute position, so
  the bond head never needs homing --- the machine knows where it is the
  moment it powers up.

\item[Non-stick] A bond that fails to weld: the wire lifts off the pad instead
  of adhering.

\item[Overtravel] How far \emph{below} the pad surface the machine aims the
  tool. Always zero or negative. It is a direction to push in rather than a
  height to arrive at --- the workpiece stops the descent.

\item[Pad] The metallised area on the workpiece that the wire is bonded to.

\item[Power] The electrical drive level for a weld, in milliwatts. Sets
  ultrasonic \emph{amplitude} --- how hard the tool scrubs.

\item[Q (quality factor)] A measure of how sharply the transducer resonates.
  Reported by \keycap{TEST}. A falling Q over time points at a loosening horn
  or a degrading stack.

\item[Reset height] The parked Z height. The machine returns here after every
  cycle and is driven here before every cycle starts.

\item[Resonance] The frequency at which the transducer converts electrical
  drive into mechanical motion most efficiently. The machine finds it by
  sweeping the scan band before every weld.

\item[Scan] The impedance sweep that locates resonance, run automatically
  before each weld across \menuitem{Scan Start} to \menuitem{Scan Stop}.

\item[Search height] The height the tool descends to quickly before beginning
  its controlled search for the pad.

\item[Stepback] The Y position the stage moves to while preparing the second
  bond.

\item[Tail] The length of wire left protruding from the tool after the wire is
  torn, ready to start the next first bond. Too short and the next bond has no
  wire to work with; too long and the loop is untidy.

\item[Tail assist] A brief ultrasonic burst applied while the tail is
  re-formed after the tear, which helps the wire feed cleanly.

\item[Tear] The action that breaks the wire after the second bond. Performed
  by the T axis in most modes, and by the Y stage in Table Tear mode.

\item[Tracking force] The force loaded while the head searches for the pad,
  before the full bond force is applied.

\item[Transducer] The ultrasonic assembly that vibrates the tool. Driven at
  its resonance by the machine.

\item[T axis] The axis that feeds the wire tail and tears the wire. It has no
  home switch, so all its moves are relative to wherever it currently sits.

\item[Wedge (tool)] The tip that presses the wire onto the pad and transmits
  ultrasonic energy into it.

\item[Y axis] The stage that carries the workpiece. Homed at power-up.

\item[Z axis] The bond head's vertical axis. Position is measured absolutely
  by the LVDT.

\end{description}
